\section{Introduction}
\label{sec:intro}

Internet memes combine an image with overlaid text, and their meaning usually arises from the interaction of the
two: a harmless caption can become sarcastic or offensive in front of a particular picture, and vice versa. This
interaction makes memes hard for systems that analyse text or images in isolation. It has motivated a line of
multimodal datasets and models for sentiment, humour, offensiveness and hate
speech~\cite{kiela2020hateful,sharma2020memotion}.

Bangla is spoken by more than 250 million people, and Bangla memes have recently received dedicated
resources~\cite{hossain2022memosen,hossain2022mute,das2023banglaabusememe,hossain2024bhm}. The most recent,
CMBan~\cite{cmban2025}, focuses on \emph{cartoon-driven} memes. These are built on recurring cartoon templates and
written in Bangla, English or a mix of both, and are annotated for humour, sarcasm, offensiveness, motivation and
sentiment.

All of these works evaluate models on held-out \emph{clean} data. Social-media memes, however, are not clean.
\begin{itemize}
  \item The same joke circulates in Bangla script and in Romanised ``Banglish'', with inconsistent spelling,
  English insertions and added or stripped emojis.
  \item Images are recompressed, cropped, resized and re-captioned as they are shared.
\end{itemize}
A model that scores well on a benchmark may therefore fail on the variants it meets in deployment, for example if
it depends on surface spellings or on image details rather than on the meme's meaning. To our knowledge, no prior
work measures how Bangla meme classifiers behave under such variation.

We present \textbf{ToonSense}, a robustness study of multimodal models on Bangla--English cartoon memes. ToonSense
addresses three questions:
\begin{enumerate}
  \item[\textbf{Q1}] How much do text, image and their fusion contribute to clean-data accuracy?
  \item[\textbf{Q2}] How much does accuracy drop under realistic linguistic and visual perturbations, and which
  modality and perturbation types are most harmful?
  \item[\textbf{Q3}] Why do the models fail? Do attribution patterns shift under perturbation, and can errors be
  traced to the text, the image, their combination or cultural context?
\end{enumerate}

Our contributions are:
\begin{itemize}
  \item \textbf{A robustness framework for Bangla memes.} It covers twelve perturbation families at three
  severities. The six linguistic families include a colloquial Bangla-to-Banglish transliterator, Romanised and
  native-script spelling noise, code-mixing and emoji edits. The six visual families include text-region
  occlusion driven by OCR boxes. Performance loss is measured by the Robustness Drop and the label-free flip rate.

  \item \textbf{A multi-encoder benchmark.} We compare six frozen multilingual text encoders, four image encoders
  and five fusion heads, using repeated cross-validation, bootstrap confidence intervals and Holm-corrected
  significance tests. Besides the publicly labelled CMBan subset, we replicate all analyses on two larger Bangla
  meme datasets, BanglaAbuseMeme~\cite{das2023banglaabusememe} and MemoSen~\cite{hossain2022memosen}.

  \item \textbf{Comparative explainability.} We compare Integrated Gradients modality shares, token attributions
  and Grad-CAM maps \emph{before and after} perturbation. We introduce a text-attention ratio that measures how
  much CLIP looks at the rendered caption.

  \item \textbf{An error taxonomy} that attributes clean errors and robustness failures to linguistic, visual,
  cross-modal, OCR-related and contextual/cultural sources.
\end{itemize}

Our main findings are as follows.
\begin{itemize}
  \item Image encoders carry most of the signal: CLIP beats every frozen text encoder on all three datasets, and a
  character $n$-gram baseline beats the text encoders on the two larger ones. With a few thousand training memes,
  fusion adds a small but significant gain.
  \item Romanised Bangla is the most damaging perturbation on every dataset. It reduces the macro-F1 of the fusion
  models by 8--21\% relative and changes 26--45\% of their predictions, while native-script spelling noise and
  code-mixing are nearly harmless.
  \item Fusion halves this damage on the larger datasets but not on the small CMBan subset. On all datasets, the
  perturbed text \emph{gains} attribution share rather than being discounted.
  \item Adding transliterated copies of the training memes removes most of the drop, including for a Romanised
  spelling style never seen in training, at a cost of under one point of clean macro-F1.
  \item Visual perturbations flip up to 27\% of predictions with comparatively small F1 changes, so aggregate
  accuracy understates instability. The largest residual error class is contextual or cultural.
\end{itemize}
Code, perturbation generators and evaluation scripts are released to support reproducible robustness evaluation of
Bangla meme models.
